\documentclass[11pt,a4paper]{article}
\usepackage[T1]{fontenc}
\usepackage[utf8]{inputenc}
\usepackage[a4paper,margin=2.2cm]{geometry}
\usepackage{amsmath}
\usepackage{hyperref}
\setlength{\parindent}{0pt}
\setlength{\parskip}{0.5em}
\title{Estimación experimental de inclinaciones con LiDAR}
\author{}
\date{}
\begin{document}
\maketitle

\section{Objetivo y datos}
Estimar inclinaciones de detecciones de paneles solares de 2023 compatibles con
una cubierta coplanar o aproximadamente paralela. Se conservan recortes LiDAR
con contexto alrededor del polígono original. Las aportaciones de varias hojas
se reúnen sin duplicados exactos. Las ortofotos se usan para revisión visual.

\section{Método}
RANSAC busca el plano interior con mayor apoyo y se reajusta por SVD. Para una
normal unitaria $\mathbf n$ y un centro $\mathbf c$, la distancia perpendicular es
$d_i=(\mathbf p_i-\mathbf c)\cdot\mathbf n$ y la inclinación es
$\beta=\arccos(|n_z|)$. Se exige un 75\,\% de apoyo interior con tolerancia
de 0,10~m, al menos diez inliers y extensión horizontal en dos direcciones.
La estabilidad se comprueba variando la tolerancia y tomando una submuestra.

En el exterior se extraen hasta seis planos por franja, retirando sucesivamente
sus inliers. Se detiene la búsqueda si no queda apoyo suficiente. Los fragmentos
se separan por proximidad y se relacionan con los inliers interiores en 3D.
Entre superficies distintas, el fragmento elegido debe tener al menos el doble
de puntos vecinos que las alternativas. Los empates se conservan como ambiguos.
Sin apoyo próximo, la franja se amplía de 2 a 4 y hasta 8~m.

La aceptación compara estabilidad, ángulo entre normales y separación: se
admiten hasta 5 grados y 0,30~m en el percentil 90. Los límites son configurables.
Se guardan por separado inclinación candidata, estimada y utilizada. Cuando
falta evidencia, la estimada queda desconocida y la utilizada toma un supuesto
de 30 grados, configurable. Esta hipótesis no equivale a una medida.

\section{Alcance y validación}
Se estima el faldón local, también en naves de poca pendiente. Varias vertientes
en una detección, desplazamientos y cubiertas sin superficie exterior visible
pueden impedir la comprobación. La continuidad geométrica no identifica por sí
sola los módulos ni confirma su montaje. Compartir año de campaña tampoco
garantiza simultaneidad entre ortofoto y LiDAR.

La validación necesita edificios distintos, incluyendo naves y cubiertas grandes,
y una referencia independiente para medir error angular. Se comparará tanto el
número de detecciones aprovechables como su superficie. Los resultados de este
prototipo todavía no sustituyen las hipótesis del cálculo energético.

En la muestra de cien casos, la búsqueda exterior anterior aceptaba trece
inclinaciones y la versión actual acepta veintiuna. Es un cambio de
aprovechamiento, no una medida de exactitud. Las salidas de las pruebas
sustituidas y del lote regional incompleto del método anterior son
prescindibles. Se conservan los recortes originales, los formularios y los
resultados actuales.

Los comandos y archivos se describen en \texttt{docs/lidar/README.md}.
Como referencia del método:
\href{https://www.open3d.org/docs/release/tutorial/geometry/pointcloud.html}{documentación de segmentación de planos de Open3D}.
\end{document}
